This chapter follows all 55 catalogue records, 174--228, through the supplied
leaves for 1969--70 and 1970--71. Source order is decisive: Hewitt's November
1969 lecture precedes the February 1970 monodromy entry, and the finite-field
Ramsey question precedes the next sheaf lecture despite its higher record
number. The opening Buchstaber fragment on P000282 belongs to the preceding
year; the Kirillov leaf on P000360 belongs to the following one. Both are
documented as boundary evidence rather than assigned to adjacent speakers
\citep[catalogue, records 173--229]{mccme2026gelfand}.

The principal continuous thread is Kazhdan's course in schemes, sheaves,
affine cohomology and projective constructions, joined to Griffiths's
periods by monodromy, residues and vanishing cycles. A second thread passes
from McDuff's factors and Sinai's quantum chains to Gelfand's current-group
questions and Perelomov's coherent states. Finite groups, four subspaces,
matrix deformations, real curves, characteristic numbers, torsion and
inverse scattering repeatedly connect these strands. Announcements and
reading programmes remain catalogue entries even when they contain no
identifiable mathematical assertion.

The occurrence index contains 290 items and the problem index
17 definite questions or conjectures. Counts measure editorial
segmentation, not originality. The historical classes are:
\begin{center}
\begin{tabular}{lrrrrrr}
\toprule
Class&A&B&C&D&E&F\\
\midrule
Items&19&8&4&0&231&28\\
\bottomrule
\end{tabular}
\end{center}
Class E dominates: the seminar mixed exposition, verification and
research questions. Chronology matters.
McDuff's continuum of factors had appeared before her October lecture;
Griffiths's period results had circulated before their 1970 journal issue;
Perelomov's lattice completeness paper preceded his April lecture.
Conversely, later general solutions must not be substituted for narrower
seminar constructions. The general modularity theorem, the Gudkov
congruence and analytic--combinatorial torsion are followed through
their later literature in their respective entries.

Two bounded developments are supplied with complete proofs and independent
checks. For current groups, an infinite tensor product of closely paired
evaluations gives a strongly continuous irreducible representation of
$C^1([0,1],SU(2))$, still irreducible on smooth currents, with support
$[0,1]$. A countersequence shows why its topology cannot be weakened to
uniform continuity. Operator-norm-continuous smooth representations have
finite support by the checked classification
\citep{janssensneeb2015norm}; the notebook's unspecified category remains
an explicit qualification. For the four-subspace question, the orbit of
$S(3,2)^{\oplus2}$ is open dense in
$\operatorname{Gr}(4,6)^4$, giving the requested deformation of
$S(5,2)\oplus S(1,2)$. Its mathematical conclusion is settled here;
the exact historical attribution of that elementary proof is not.
Neither development claims novelty.

Two further annotations audit settled material. A faithful Bell-block
state disproves the claim that commuting finite marginals necessarily
have compatible one-site product eigenbases. Exact rational computations
check all four-site marginals and commutators at two parameter values.
Wavelet scaling verifies the positive answer to the product-uncertainty
question, without validating the separate proposed Gram--Schmidt method.
The arbitrary-sequence coherent-state completeness problem remains
qualified: checked partial criteria do not settle all irregular sets,
and an inaccessible September 2026 preprint prevents a certified global
openness or resolution claim as of 9 November 2026.

The 99 corrections distinguish transcription, translation,
ambiguity, provenance and faithful mathematical error. Earlier-year links
are retained where supported; related themes are not merged into one
problem merely because their titles resemble one another. In particular,
the earlier distribution/foliation slogan is a contextual link, while
the Buchstaber carryover keeps its existing identity. Throughout, the
notes, modern reconstruction, publication match and present mathematical
status are separately visible.
